\documentclass[11pt]{article}
% ===========================================================================
%  Shared preamble for the AP Calculus AB exam-preparation worksheets.
%  Input right after \documentclass. Files beginning with "_" are not compiled
%  on their own. Navy and gold, to match the course's lesson pages.
% ===========================================================================
\usepackage[letterpaper,top=2.3cm,bottom=2.2cm,left=1.7cm,right=1.7cm,headheight=28pt,footskip=30pt]{geometry}
\usepackage{amsmath,amssymb}
\usepackage[T1]{fontenc}
\usepackage[utf8]{inputenc}
\usepackage{lmodern}
\usepackage{xcolor}
\usepackage[most]{tcolorbox}
\usepackage{enumitem}
\usepackage{fancyhdr}
\usepackage{array}
\usepackage{tikz}
\usepackage{pgfplots}
\pgfplotsset{compat=1.18}

\definecolor{navy}{HTML}{161E54}
\definecolor{ink}{HTML}{27357E}
\definecolor{gold}{HTML}{B98A00}
\definecolor{goldbg}{HTML}{FFFDF3}
\definecolor{inkbg}{HTML}{F6F7FD}
\definecolor{rule}{HTML}{C9CFF0}
\definecolor{okgreen}{HTML}{1B7A44}
\definecolor{okbg}{HTML}{F1F8F3}

\pagestyle{fancy}
\fancyhf{}
\fancyhead[L]{\footnotesize Name: \underline{\hspace{4cm}}}
\fancyhead[C]{\textbf{\color{navy}Integration Academy}\\[-2pt]{\scriptsize Calculus AB \textperiodcentered{} AP\textsuperscript{\textregistered} Exam Preparation}}
\fancyfoot[L]{\tiny\color{black!55}\parbox[t]{0.78\textwidth}{AP\textsuperscript{\textregistered} is a registered trademark of the College Board, which was not involved in the production of, and does not endorse, this material. \textcopyright{} Integration Academy.}}
\fancyfoot[R]{\scriptsize\thepage}
\renewcommand{\headrulewidth}{1.2pt}
\renewcommand{\headrule}{\hbox to\headwidth{\color{gold}\leaders\hrule height \headrulewidth\hfill}}
\renewcommand{\footrulewidth}{0pt}
\setlength{\parindent}{0pt}
\setlength{\parskip}{4pt}

% Title block: \wstitle{code}{title}{unit line}
\newcommand{\wstitle}[3]{\begin{center}{\small\bfseries\color{gold}WORKSHEET #1}\\[3pt]{\LARGE\bfseries\color{navy}#2}\\[3pt]{\small\color{black!60}#3}\end{center}\vspace{2pt}}
% Section heading
\newcommand{\sech}[1]{\par\vspace{10pt}{\large\bfseries\color{navy}#1}\par\nobreak\vspace{2pt}{\color{rule}\hrule height0.8pt}\vspace{7pt}}
% A level tag printed beside a group of questions
\newcommand{\level}[1]{\par\vspace{6pt}{\footnotesize\bfseries\color{gold}\MakeUppercase{#1}}\par\nobreak\vspace{2pt}}

\newtcolorbox{ideabox}{colframe=ink,colback=inkbg,coltitle=white,colbacktitle=ink,fonttitle=\bfseries\small,title={Key ideas},arc=2pt,boxrule=0.8pt,left=8pt,right=8pt}
\newtcolorbox{example}[1]{colframe=rule,colback=white,coltitle=navy,colbacktitle=inkbg,fonttitle=\bfseries\small,title={#1},arc=2pt,boxrule=0.8pt,borderline west={3pt}{0pt}{ink},left=10pt,right=8pt}
\newtcolorbox{frbox}[1]{colframe=gold,colback=goldbg,coltitle=white,colbacktitle=gold,fonttitle=\bfseries\small,title={#1},arc=2pt,boxrule=0.8pt,breakable,left=8pt,right=8pt}
\newtcolorbox{solbox}{colframe=okgreen,colback=okbg,coltitle=white,colbacktitle=okgreen,fonttitle=\bfseries\small,title={Worked solutions},arc=2pt,boxrule=0.8pt,breakable,left=8pt,right=8pt}
\newtcolorbox{scorebox}[1]{colframe=navy,colback=white,coltitle=white,colbacktitle=navy,fonttitle=\bfseries\small,title={#1},arc=2pt,boxrule=0.8pt,breakable,left=8pt,right=8pt}

\newcommand{\soln}{\par\smallskip\textbf{\color{ink}Solution.}\ }
\newcommand{\result}[1]{\par\smallskip\textbf{#1}}
% blank working space of a given height
\newcommand{\work}[1]{\par\vspace{#1}}
% a point value in the margin of a scoring guide
\newcommand{\pts}[1]{\hfill{\footnotesize\bfseries\color{gold}[#1]}}

% Coordinate grid for graph questions: \begin{apgraph}{xmin}{xmax}{ymin}{ymax} ... \end{apgraph}
\newenvironment{apgraph}[4]{\begin{center}\begin{tikzpicture}\begin{axis}[
  axis lines=middle, xmin=#1, xmax=#2, ymin=#3, ymax=#4,
  grid=both, grid style={black!12}, xtick distance=1, ytick distance=1,
  ticklabel style={font=\scriptsize}, xlabel={$x$}, ylabel={$y$},
  width=10cm, height=7.2cm, samples=80, clip=false,
  every axis plot/.append style={line width=1.1pt, color=ink}]}{\end{axis}\end{tikzpicture}\end{center}}
\newcommand{\closeddot}[2]{\addplot[only marks, mark=*, mark size=2.2pt, color=ink] coordinates {(#1,#2)};}
\newcommand{\opendot}[2]{\addplot[only marks, mark=*, mark size=2.2pt, color=ink, mark options={fill=white}] coordinates {(#1,#2)};}

% Multiple-choice options: four across, or stacked when the options are long.
\newcommand{\opts}[4]{\par\vspace{3pt}\begin{tabular}{@{}p{0.235\linewidth}p{0.235\linewidth}p{0.235\linewidth}p{0.235\linewidth}@{}}(A)\ #1 & (B)\ #2 & (C)\ #3 & (D)\ #4\end{tabular}\par\vspace{9pt}}
\newcommand{\optsv}[4]{\par\vspace{3pt}\begin{tabular}{@{}l@{\ \ }p{0.9\linewidth}@{}}(A) & #1\\(B) & #2\\(C) & #3\\(D) & #4\end{tabular}\par\vspace{9pt}}

\fancyhead[R]{\footnotesize\color{black!60}Worksheet 6.7}
\begin{document}

\wstitle{6.7}{The Fundamental Theorem, Part 2: Evaluating Integrals}{Unit 6: Integration and Accumulation of Change}

\begin{ideabox}
\begin{itemize}[leftmargin=*,itemsep=2pt,topsep=2pt]
  \item \textbf{Fundamental Theorem, Part 2.} If $F'=f$, then $\displaystyle\int_{a}^{b}f(x)\,dx=F(b)-F(a)$.
  \item Write the antiderivative in brackets with the limits, then substitute: upper limit first.
  \item $\displaystyle\int_{a}^{b}f'(x)\,dx=f(b)-f(a)$: the integral of a rate of change is the net change. So $f(b)=f(a)+\displaystyle\int_{a}^{b}f'(x)\,dx$.
  \item For an absolute value or a piecewise function, split the interval where the formula changes.
\end{itemize}
\end{ideabox}

\sech{Worked examples}
\begin{example}{Example 1 \textperiodcentered{} A polynomial}
Evaluate $\displaystyle\int_{0}^{2}\left(3x^2+2x\right)\,dx$.\soln $\Big[x^3+x^2\Big]_0^2=(8+4)-0$.\result{$12$.}
\end{example}

\begin{example}{Example 2 \textperiodcentered{} A trigonometric function}
Evaluate $\displaystyle\int_{\pi/6}^{\pi/2}\cos x\,dx$.\soln $\Big[\sin x\Big]_{\pi/6}^{\pi/2}=1-\dfrac12$.\result{$\dfrac12$.}
\end{example}

\begin{example}{Example 3 \textperiodcentered{} A value from a rate}
Given $f(1)=4$ and $f'(x)=2x+3$, find $f(3)$.\soln $f(3)=4+\displaystyle\int_{1}^{3}(2x+3)\,dx=4+\Big[x^2+3x\Big]_1^3=4+(18-4)$.\result{$18$.}
\end{example}

\newpage
\sech{Practice}

\level{Core}
\begin{enumerate}[leftmargin=*,itemsep=2pt]
  \item Evaluate $\displaystyle\int_{1}^{3}2x\,dx$.\work{2.3cm}
  \item Evaluate $\displaystyle\int_{1}^{4}\dfrac{1}{\sqrt{x}}\,dx$.\work{2.3cm}
  \item Evaluate $\displaystyle\int_{0}^{\ln2}e^{x}\,dx$.\work{2.3cm}
  \item Evaluate $\displaystyle\int_{1}^{2}\dfrac{3}{x^2}\,dx$.\work{2.3cm}
\end{enumerate}

\level{Exam level}
\begin{enumerate}[leftmargin=*,itemsep=2pt,start=5]
  \item Evaluate $\displaystyle\int_{1}^{8}\sqrt[3]{x}\,dx$.\work{2.8cm}
  \item Evaluate $\displaystyle\int_{0}^{\pi/3}\sec x\tan x\,dx$.\work{2.8cm}
  \item Let $f(x)=x^2$ for $x\le2$ and $f(x)=4$ for $x>2$. Evaluate $\displaystyle\int_{0}^{5}f(x)\,dx$.\work{2.8cm}
  \item A particle has velocity $v(t)=3t^2-6t$ and position $x(0)=2$. Find $x(3)$.\work{2.8cm}
\end{enumerate}

\begin{frbox}{Question 9 \textperiodcentered{} Free-response practice}
Answer each part. No calculator is permitted.
\begin{enumerate}[label=(\alph*),leftmargin=*,itemsep=2pt]
  \item Evaluate $\displaystyle\int_{1}^{e}\left(\dfrac2x+1\right)\,dx$.\work{2.6cm}
  \item A function $f$ has $f(0)=5$ and $f'(x)=4\cos x$. Find $f\!\left(\dfrac\pi2\right)$.\work{2.6cm}
  \item Find the area of the region under $y=4-x^2$ and above the $x$-axis for $0\le x\le2$.\work{2.6cm}
\end{enumerate}
\end{frbox}

\newpage
\begin{solbox}
\begin{enumerate}[leftmargin=*,itemsep=4pt]
  \item $\Big[x^2\Big]_1^3=8$.
  \item $\Big[2\sqrt x\Big]_1^4=4-2=2$.
  \item $e^{\ln2}-e^0=1$.
  \item $\Big[-\dfrac3x\Big]_1^2=-\dfrac32+3=\dfrac32$.
  \item $\Big[\dfrac34x^{4/3}\Big]_1^8=\dfrac34(16-1)=\dfrac{45}{4}$.
  \item $\Big[\sec x\Big]_0^{\pi/3}=2-1=1$.
  \item $\displaystyle\int_{0}^{2}x^2\,dx+\displaystyle\int_{2}^{5}4\,dx=\dfrac83+12=\dfrac{44}{3}$.
  \item $2+\Big[t^3-3t^2\Big]_0^3=2+0=2$.
  \item \begin{enumerate}[label=(\alph*),leftmargin=*,itemsep=3pt]
    \item $\Big[2\ln x+x\Big]_1^e=(2+e)-(0+1)=e+1$.
    \item $5+\displaystyle\int_{0}^{\pi/2}4\cos x\,dx=5+4=9$.
    \item $\Big[4x-\dfrac{x^3}{3}\Big]_0^2=8-\dfrac83=\dfrac{16}{3}$.
  \end{enumerate}
\end{enumerate}
\end{solbox}

\end{document}
